\chapter{Architecture and Design Principles}\label{ch:arch}

\section{Goals and non-goals}

The design goals, as recorded in the repository's roadmap~\cite{repo}:
\begin{enumerate}
  \item Per-stream reliability, ordering and congestion control.
  \item High throughput on commodity hosts, through batching and the absence of per-packet allocation.
  \item Optional, fail-closed encryption: a peer that is configured to encrypt never silently falls back to
    plaintext.
  \item Robustness against hostile input: every field received from the network is treated as an attacker's
    claim.
  \item No workload-specific assumptions in the protocol core.
\end{enumerate}

\noindent Things the protocol deliberately does not do:
\begin{itemize}
  \item renegotiate block or window sizes in the middle of a session;
  \item make encryption mandatory;
  \item authenticate peers with certificates or a PKI (pre-shared keys are the design);
  \item expose a byte-stream socket compatible with TCP.
\end{itemize}

\section{Layered structure}

Figure~\ref{fig:layers} shows the structure of the repository:
\begin{itemize}
  \item \textbf{Protocol core} (\code{include/fuse/proto/}, \code{src/proto/}): small building blocks, each
    independently testable and none of them aware of sockets, threads or clocks.
  \item \textbf{Transports}: built from those blocks, and the only place where the blocks meet real I/O,
    real time and real concurrency.
  \item \textbf{UDP wrapper}: the lowest layer, which hides the platform differences.
\end{itemize}

\begin{figure}[H]
\centering
\fitwidth{\begin{tikzpicture}[node distance=4mm and 3mm, every node/.style={font=\small}]
  \node[api,minimum width=3.3cm] (t) {\textbf{transfer.hpp}\\buffers \& files};
  \node[api,minimum width=3.3cm,right=of t] (s) {\textbf{sdk.h}\\connect / send / recv};
  \node[api,minimum width=3.3cm,right=of s] (m) {\textbf{mux\_transfer.hpp}\\N streams, 1 socket};
  \node[api,minimum width=3.0cm,right=of m] (h) {\textbf{http.h}\\1 request, 1 response};
  \node[core,below=10mm of t,minimum width=2.6cm] (w) {wire / block / control\\encode \& decode};
  \node[core,right=of w,minimum width=2.4cm] (r) {registry\\retransmit store};
  \node[core,right=of r,minimum width=2.4cm] (rx) {receiver\\window, NACK};
  \node[core,right=of rx,minimum width=2.6cm] (cc) {congestion\\AIMD, RTO, pacing};
  \node[core,right=of cc,minimum width=2.2cm] (cr) {session\_crypto\\AEAD, GMAC};
  \node[core,below=4mm of w,minimum width=2.6cm] (su) {setup\\3-way handshake};
  \node[core,right=of su,minimum width=2.4cm] (q) {send\_queue,\\reassembly};
  \node[core,right=of q,minimum width=2.4cm] (ring) {spsc\_ring\\lock-free};
  \node[core,right=of ring,minimum width=2.6cm] (o) {orchestrator\\autoscaling};
  \node[core,right=of o,minimum width=2.2cm] (d) {digest, hash,\\dtls};
  \node[io,below=10mm of ring,minimum width=9cm] (u) {\textbf{udp.hpp}: GSO + \code{recvmmsg} (Linux), plain sockets (macOS)};
  \node[io,below=4mm of u,minimum width=9cm] (k) {Kernel UDP / IP};
  \begin{scope}[on background layer]
    \node[draw=black!30,dashed,rounded corners,fit=(w)(cr)(su)(d),inner sep=6pt] {};
  \end{scope}
  \draw[arr] (h) -- (s);
  \draw[arr] ($(t.south)+(0,0)$) -- ++(0,-5mm);
  \draw[arr] ($(s.south)+(0,0)$) -- ++(0,-5mm);
  \draw[arr] ($(m.south)+(0,0)$) -- ++(0,-5mm);
  \draw[arr] (u) -- (k);
  \draw[arr] ($(q.south)!0.5!(ring.south)+(0,-2mm)$) -- (u.north);
\end{tikzpicture}}
\caption{Layered structure of the repository: public APIs (green), protocol core (blue), platform I/O (orange).}\label{fig:layers}
\end{figure}

The public APIs are:
\begin{itemize}
  \item \code{transfer.hpp}: \code{send_buffer}, \code{receive_buffer}, \code{send_file} and \code{receive_file};
  \item \code{sdk.h}: a C API with \code{fuse_listen}, \code{fuse_accept}, \code{fuse_connect},
    \code{fuse_send}, \code{fuse_recv} and \code{fuse_close};
  \item \code{mux_transfer.hpp}: \code{send_multiplexed} and \code{receive_multiplexed};
  \item \code{http.h}: one HTTP/1.1 request and one response over an SDK connection.
\end{itemize}

They build into the libraries listed in Table~\ref{tab:targets}.

\begin{table}[H]
\centering\small
\caption{CMake library targets.}\label{tab:targets}
\begin{tabularx}{\linewidth}{lXl}
\toprule
\textbf{Target} & \textbf{Contents} & \textbf{Installed} \\
\midrule
\code{fuse::proto} & Protocol core: block, control, registry, receiver, UDP, SETUP, reassembly, DTLS, orchestrator, session crypto, digest & yes \\
\code{fuse::sdk} & The three transports (\code{transfer.cpp}, \code{mux_transfer.cpp}, \code{sdk.cpp}), linking \code{fuse::proto} & no$^{*}$ \\
\code{fuse::http} & HTTP mapping (\code{http.cpp}) & no$^{*}$ \\
\code{fuse::fuse} & Legacy C scaffold (QUIC-like long header, variable-length integers); unused by the protocol & yes \\
\bottomrule
\multicolumn{3}{l}{\footnotesize $^{*}$Released from the separate \code{fuse-sdk} repository.}
\end{tabularx}
\end{table}

Two build options matter most:
\begin{itemize}
  \item \code{FUSE_WITH_CRYPTO} (on by default) links wolfSSL~\cite{wolfssl}. The build can fetch wolfSSL and
    compiles it with AES-NI assembly: wolfSSL's own CMake build had silently used software AES, at 151~MB/s
    against 4,852~MB/s with the assembly enabled.
  \item Without crypto, every cryptographic entry point fails closed, and a request for encryption returns
    \code{Unsupported}.
\end{itemize}

\section{The three transports at a glance}

\begin{table}[H]
\centering\small
\caption{The three transports.}\label{tab:transports}
\begin{tabularx}{\linewidth}{>{\raggedright\arraybackslash}p{2.4cm}>{\raggedright\arraybackslash}X>{\raggedright\arraybackslash}X>{\raggedright\arraybackslash}X}
\toprule
 & \textbf{Bulk transfer} & \textbf{SDK} & \textbf{Multiplexed} \\
\midrule
Unit & one buffer or file, split into \emph{lanes} & messages on a connection & one buffer over $N$ streams \\
Ports & one per lane (\code{base_port}$+i$) & one listen port, then a fresh port per connection & one \\
Threads & sender: one per lane; receiver: one reader per lane, $K$ writers, one retransmission consumer & one pump thread per connection & autoscaling pool \\
Session setup & per-lane StreamStart with Retry token\vviii & HELLO / HELLO-ACK & SETUP, then per-stream StreamStart \\
Encryption & AES-256-GCM per block; GMAC on control\vviii; key epochs\vviii & AES-256-GCM on DATA only & none \\
v8 hardening & complete & deferred & SACK, RTO, slow start, validation \\
\bottomrule
\end{tabularx}
\end{table}

\section{Evolution of the wire protocol}

The first byte of every datagram is the protocol version. Each increment was an on-wire change, and the
comment above the constant in \code{wire.hpp} records why (Table~\ref{tab:versions}).

\begin{table}[H]
\centering\small
\caption{Protocol versions and the reason for each change.}\label{tab:versions}
\begin{tabularx}{\linewidth}{c>{\raggedright\arraybackslash}p{5cm}X}
\toprule
\textbf{v} & \textbf{Change} & \textbf{Reason} \\
\midrule
1 & Block header, registry, bitmask, ACK/NACK & Stage 1 reliability core \\
2 & Explicit byte offset in DATA & Block size adapts mid-stream, so $\mathit{seq}\times\mathit{size}$ misplaces data \\
3 & Multi-word bitmask; handshake nonce & Windows over 64 blocks; a stale ACK must not confirm a new session \\
4 & Server salt in HELLO-ACK & Retried HELLOs could derive identical keys with both sequence numbers at 0, an AES-GCM nonce collision \\
5 & Two-byte window size & Raise the window from 255 to 1,024 blocks; the old field capped throughput on any real bandwidth-delay product \\
6 & Base offset and total size in StreamStart & Write each block directly to its final file position \\
7 & File id, ResumeRanges, FileDigest & Resumable file reception with whole-file verification \\
8 & Control-plane authentication, AAD fix, key epochs, Retry, path validation, SETUP nonce, receive window, extensions, version-mismatch reply & The security and correctness review (Chapter~\ref{ch:security}) \\
\bottomrule
\end{tabularx}
\end{table}

\section{Design principles of the implementation}\label{sec:principles}

The code base is small, but it follows a consistent set of principles. They explain many choices that look
unusual in isolation.

\paragraph{Staged construction with acceptance tests.} The protocol was built in stages, each with its own
acceptance test:

\smallskip
\begin{center}\small
\begin{tabular}{@{}ll@{\qquad}ll@{}}
\toprule
\textbf{Stage} & \textbf{Scope} & \textbf{Stage} & \textbf{Scope} \\
\midrule
0 & framing & 4 & per-stream flags \\
1 & blocks and the received bitmask & 5 & congestion control \\
1.5 & control messages & 6 & CPU pinning \\
2 & SETUP & 7 & encryption \\
3 & worker threads & 8 & benchmarks \\
\bottomrule
\end{tabular}
\end{center}
\smallskip

\noindent When a measurement contradicted a stage, the stage was removed rather than kept because it had been
planned. CPU pinning made transfers about 6\% slower, so Stage~6 was deleted, and the reason is recorded in
\code{orchestrator.hpp}.

\paragraph{Pure state machines; the caller owns I/O and time.} The SETUP handshake, the receiver window, the
sender registry and the congestion controller never touch a socket or read a clock. They receive the current
time as an argument and write any datagram they want to send into a buffer supplied by the caller. As a result:
\begin{itemize}
  \item tests can drive them with a virtual clock and injected loss, so ``retransmits after 50~ms'' is a
    deterministic unit test rather than a sleep;
  \item the same logic serves three transports with different threading models.
\end{itemize}
Loops that depend on time are bounded: the congestion controller processes at most 4,096 epochs per call, so a
clock jump cannot make it spin.

\paragraph{A single place for byte order.} Only the \code{put_u8}\ldots\code{put_u64} and
\code{get_u8}\ldots\code{get_u64} helpers in \code{wire.hpp} touch raw wire bytes. They are all big-endian, and
each returns the number of bytes it consumed, so cursor arithmetic looks the same everywhere. Every message has
an encoder that returns the encoded length (zero if it does not fit) and a decoder that returns false on a
version mismatch, the wrong type or truncation. A malformed datagram is therefore dropped at one obvious place.

\paragraph{Fixed-size storage and no allocation on the hot path.}
\begin{itemize}
  \item A registry slot embeds its whole payload array rather than a pointer.
  \item The receiver's bitmask and gap timers are fixed arrays sized by the maximum window.
  \item Messages are plain structures with fixed arrays.
\end{itemize}
A test that replaces the global \code{operator new} proves that the Stage~1 send and receive loop allocates
nothing after setup. This avoids allocator latency and jitter, and a hostile length field cannot grow a fixed
array. The cost is memory sized for the worst case, which is visible in the SDK (Section~\ref{sec:sdk}).

\paragraph{Every wire field is hostile.} Lengths, offsets and counts from the network are bounded before any
arithmetic is done with them, and in an order that cannot overflow. For example, a block's placement is checked
as $\mathit{offset}\le\mathit{total} \wedge \mathit{len}\le\mathit{total}-\mathit{offset}$, never as
$\mathit{offset}+\mathit{len}\le\mathit{total}$, because an offset close to $2^{64}$ would wrap the sum. A claimed
transfer size that cannot be allocated fails the transfer with \code{ResourceLimit} instead of throwing out of a
thread.

\paragraph{Fail closed.}
\begin{itemize}
  \item With no cryptographic backend, encryption requests are refused rather than downgraded.
  \item With no secure random source, the process aborts rather than use a predictable nonce.
  \item If sealing a block fails, the batch stops rather than emit plaintext.
  \item A cipher object refuses to be keyed twice, because re-keying is how a (key, nonce) pair could repeat.
  \item Since version~8, under a pre-shared key, nothing unsigned is acted upon.
\end{itemize}

\paragraph{Ownership through RAII.}
\begin{itemize}
  \item Sockets are move-only objects, and ciphers are uniquely owned.
  \item The digest thread of a sending lane is joined on every exit path by a small guard object.
  \item The SDK counts the application threads inside \code{fuse_send} and \code{fuse_recv} with an RAII guard,
    so that closing a connection can wait for them.
  \item Key material is wiped after use.
\end{itemize}
Lane functions have a dozen early exits (cancellation, timeout, abort, authentication failure), and RAII is what
keeps the cleanup correct on all of them.

\paragraph{A concurrency model chosen per component.} Three patterns are used, each where it fits:
\begin{enumerate}
  \item \emph{Thread confinement}. A bulk-transfer lane owns its socket, registry, receiver state and ciphers
    outright, so nothing on its per-packet path is locked.
  \item \emph{Lock-free single-producer, single-consumer rings} between a lane's socket reader and the writer
    threads.
  \item \emph{One mutex per stream or connection} in the multiplexed transfer and the SDK, where several
    threads touch the same stream.
\end{enumerate}
Shared counters use the weakest correct memory ordering: relaxed for progress statistics, release and acquire
for the flags that publish state.

\paragraph{Hiding dependencies.}
\begin{itemize}
  \item wolfSSL types appear in no public header; a cipher holds its context as an opaque pointer.
  \item Platform differences are confined to three files: \code{udp_linux.cpp} (segmentation offload,
    \code{recvmmsg}, MTU query), \code{udp_darwin.cpp} (one system call per datagram) and
    \code{udp_common.cpp}.
\end{itemize}

\paragraph{Time, not counts.}
\begin{itemize}
  \item All timing uses a monotonic nanosecond clock, and differences are taken modulo $2^{64}$.
  \item Retries are bounded by elapsed time. The original StreamStart handshake gave up after
    $200\times200~\mu$s $=40$~ms, so it failed on any path with a longer round trip. It now retries with
    exponential backoff until a connect timeout.
  \item The sender keeps two progress clocks. One paces retransmission. The other, which only moves when an
    acknowledgement confirms new data, enforces the user's timeout. A local send to a dead peer keeps
    ``succeeding'', so using the first clock for the timeout made the loop unbounded.
\end{itemize}

\paragraph{A C ABI where stability matters.} The SDK is plain C. It uses opaque handles, a status enumeration,
and a configuration structure initialised by a function, so that fields can be added without breaking callers.
It is the stable interface for bindings to other languages.

\paragraph{Comments state the invariant and its history.} Comments explain \emph{why}, often naming the bug
that motivated a line and the roadmap number of its fix. Protocol code is mostly invariants, and a future change
that breaks one should run into the comment that states it.
